\begin{table}[t]
\centering
\tiny
\caption{AIOpsStressBench benchmark card.}
\label{tab:benchmark-card}
\resizebox{\columnwidth}{!}{%
\begin{tabular}{p{0.18\linewidth}p{0.32\linewidth}p{0.39\linewidth}}
\toprule
Component & Setting & Details \\
\midrule
Task & Operational forecasting under deployment stress & Input length 96, horizon 24, chronological 65/15/20 split. \\
Public sources & GAIA, NetMan, Alibaba 2018, Salesforce/Borg 2011 & Alibaba and Salesforce/Borg are multivariate telemetry; GAIA/NetMan provide KPI stress diversity. \\
Stress operators & Missing points, metric outage, delayed tail, noise, burst, level shift & Stress is injected only into the input window; targets remain clean. \\
Metrics & MSE/MAE plus deployment metrics & P95 latency, parameters, memory, severity slope, and forecast-to-capacity proxy. \\
Core comparable pool & DLinear, RACE-DLinear, PatchTST-lite, official PatchTST, native iTransformer, TimeMixer & Native or native-wrapper stress protocol where available; RACE-DLinear is a lightweight stress-aware baseline used for robustness analysis. \\
Reference-only pool & LastValue, Chronos-Bolt zero-shot, LTSF-bridge iTransformer & Context, lower-bound, and feasibility references; these rows do not drive the central model-ranking claims. \\
Reproducibility & Minimal and extended reproduction paths & Minimal path regenerates the core winner and latency tables; extended path regenerates severity, sensitivity, and multi-seed results. \\
\bottomrule
\end{tabular}
}
\end{table}
